\documentclass[11pt,a4paper]{article}
\usepackage[margin=0.85in]{geometry}
\usepackage{amsmath,amssymb,mathtools}
\usepackage{enumitem,microtype,xcolor,hyperref}
\usepackage[most]{tcolorbox}
\hypersetup{colorlinks=true,linkcolor=blue!55!black}
\setlist{itemsep=.3em,topsep=.35em}
\newcommand{\N}{\mathbb N}
\newcommand{\Z}{\mathbb Z}
\newcommand{\Q}{\mathbb Q}
\newcommand{\R}{\mathbb R}
\newcommand{\C}{\mathbb C}
\newcommand{\PP}{\mathcal P}
\newcommand{\id}{\operatorname{id}}
\newcommand{\angles}[1]{\langle #1\rangle}
\newtcolorbox{examquestion}[1]{breakable,colback=blue!3,colframe=blue!50!black,title={#1}}
\newtcolorbox{worked}[1]{breakable,colback=green!3,colframe=green!45!black,title={#1}}

\begin{document}
\begin{center}
{\LARGE\bfseries Algebra I: Midsemester Practice Paper}\\[.4em]
{\large One-hour, Quiz-1-style mock}\quad\textbullet\quad 20 marks
\end{center}

\noindent\textbf{Instructions.} Allow exactly 60 minutes. Attempt \emph{one} option from each of Questions 1--4. Question 5 is compulsory. Write one reason or counterexample for every True/False answer. No notes during the attempt. The complete solutions begin after a page break; do not look at them until the hour is over.

\noindent\textbf{Suggested timing.} Questions 1 and 2: 8 minutes each; Questions 3 and 4: 12 minutes each; Question 5: 15 minutes; checking: 5 minutes.

\begin{examquestion}{Question 1 (3 marks): choose (a) or (b)}
\textbf{(a)} Let $f:X\to Y$, $A\subseteq X$, and $B\subseteq Y$. Define
\[
f_*(A)=\{f(x):x\in A\},\qquad
f^*(B)=\{x\in X:f(x)\in B\}.
\]
Prove that $f_*(A)\subseteq B$ if and only if $A\subseteq f^*(B)$. Deduce that $A\subseteq f^*f_*(A)$.

\medskip\noindent\textbf{OR}

\medskip\noindent\textbf{(b)} Let $X=\{1,2,3\}$, $Y=\{a,b,c\}$, and let $f:X\to Y$ satisfy $f(1)=f(2)=a$, $f(3)=b$. For $A=\{1,3\}$, calculate $f_*(A)$ and
\[
f_!(A)=\{y\in Y:f^{-1}(\{y\})\subseteq A\}.
\]
Explain why $c$ belongs to $f_!(A)$ even though $c\notin f(X)$.
\end{examquestion}

\begin{examquestion}{Question 2 (3 marks): choose (a) or (b)}
\textbf{(a)} Let $G,H$ be groups and let $\phi:G\to H$ satisfy $\phi(xy)=\phi(x)\phi(y)$ for every $x,y\in G$. Prove that $\phi(e_G)=e_H$ and $\phi(x^{-1})=\phi(x)^{-1}$.

\medskip\noindent\textbf{OR}

\medskip\noindent\textbf{(b)} On $S=\R\setminus\{-1\}$ define $a*b=a+b+ab$. Prove that $(S,*)$ is a group. State its identity and the inverse of $a$.
\end{examquestion}

\begin{examquestion}{Question 3 (4 marks): choose (a) or (b)}
\textbf{(a)} Consider
\[
\phi:\Z/12\Z\longrightarrow\Z/8\Z,
\qquad \phi([k]_{12})=[2k]_8.
\]
Prove that $\phi$ is well-defined and is a group homomorphism. Find its kernel and image.

\medskip\noindent\textbf{OR}

\medskip\noindent\textbf{(b)} List every element of $\angles{[7]}$ in $\Z/18\Z$ and of $\angles{[5]}$ in $\Z/20\Z$. Give a short reason for the number of elements in each subgroup.
\end{examquestion}

\begin{examquestion}{Question 4 (4 marks): choose (a) or (b)}
\textbf{(a)} Let
\[
P=\langle x,y,z\mid yxyz^{-2}=e\rangle.
\]
Show that $P$ is isomorphic to the free group on $y,z$. State how the generator $x$ is expressed using $y,z$.

\medskip\noindent\textbf{OR}

\medskip\noindent\textbf{(b)} Let $G$ be a group and $X$ a nonempty set. On the functions $G^X=\{f:X\to G\}$ define $(fg)(x)=f(x)g(x)$. Prove that $G^X$ is a group by giving its identity and the inverse of an arbitrary $f$, and checking associativity.
\end{examquestion}

\begin{examquestion}{Question 5 (compulsory; 6 marks)}
Mark each statement \emph{True} or \emph{False}. Give one sentence of justification or one counterexample.
\begin{enumerate}[label=(\alph*)]
\item For every $f:X\to Y$ and $B,C\subseteq Y$, $f^*(B\cup C)=f^*(B)\cup f^*(C)$.
\item For every $f:X\to Y$ and $A,D\subseteq X$, $f_*(A\cap D)=f_*(A)\cap f_*(D)$.
\item If $y\notin f(X)$, then $y\in f_!(A)$ for every $A\subseteq X$.
\item $\angles{[3]}=\Z/16\Z$ under addition modulo $16$.
\item There is a nontrivial group homomorphism $\Z/7\Z\to\Z/15\Z$.
\item If $g^2=e$ for every element $g$ of a group $G$, then $G$ is abelian.
\end{enumerate}
\end{examquestion}

\clearpage
\begin{center}
{\Large\bfseries Worked solutions and marking guide}\\[.3em]
\textit{Read this only after completing the timed attempt.}
\end{center}

\begin{worked}{Question 1(a)}
Suppose $f_*(A)\subseteq B$. Take any $x\in A$. Then $f(x)\in f_*(A)\subseteq B$, so $x\in f^*(B)$. Hence $A\subseteq f^*(B)$.

Conversely, suppose $A\subseteq f^*(B)$. Take any $y\in f_*(A)$. By definition $y=f(x)$ for some $x\in A$. Since $x\in f^*(B)$, we have $f(x)=y\in B$. Hence $f_*(A)\subseteq B$.

Now take $B=f_*(A)$. The inclusion $f_*(A)\subseteq f_*(A)$ is automatic, so the equivalence just proved gives $A\subseteq f^*f_*(A)$.
\end{worked}

\begin{worked}{Question 1(b)}
$A=\{1,3\}$, so $f_*(A)=\{f(1),f(3)\}=\{a,b\}$. The fibres are
\[
f^{-1}(\{a\})=\{1,2\},\quad f^{-1}(\{b\})=\{3\},\quad f^{-1}(\{c\})=\varnothing.
\]
The first fibre is not contained in $A$ because $2\notin A$. The second fibre is contained in $A$, and the empty fibre is contained in every set. Thus $f_!(A)=\{b,c\}$. In particular $c\in f_!(A)$ because its fibre is empty, so the condition ``every preimage of $c$ lies in $A$'' is true vacuously.
\end{worked}

\begin{worked}{Question 2(a)}
Let $h=\phi(e_G)$. Since $e_Ge_G=e_G$ and $\phi$ preserves products,
\[
h=\phi(e_G)=\phi(e_Ge_G)=\phi(e_G)\phi(e_G)=h^2.
\]
Multiply by $h^{-1}$ in $H$ to obtain $e_H=h$, so $\phi(e_G)=e_H$.

For any $x\in G$,
\[
\phi(x)\phi(x^{-1})=\phi(xx^{-1})=\phi(e_G)=e_H.
\]
Similarly $\phi(x^{-1})\phi(x)=e_H$. Thus $\phi(x^{-1})$ is the two-sided inverse of $\phi(x)$, and uniqueness of inverses gives $\phi(x^{-1})=\phi(x)^{-1}$.
\end{worked}

\begin{worked}{Question 2(b)}
Observe the useful identity
\[
1+(a*b)=1+a+b+ab=(1+a)(1+b).
\]
Since $a,b\ne-1$, both factors on the right are nonzero; hence $1+(a*b)\ne0$ and $a*b\ne-1$. This proves closure in $S$.

The map $T:S\to\R\setminus\{0\}$, $T(a)=1+a$, is bijective and satisfies $T(a*b)=T(a)T(b)$. Ordinary multiplication of nonzero reals is associative, so $*$ is associative. The identity is $0$ because $a*0=0*a=a$. To find an inverse $b$, solve
\[
0=a*b=a+b+ab
\quad\Longleftrightarrow\quad
b(1+a)=-a
\quad\Longleftrightarrow\quad
b=-\frac{a}{1+a}.
\]
This lies in $S$, since $1+b=1/(1+a)\ne0$. The operation is commutative, so the same $b$ is a two-sided inverse. Thus $(S,*)$ is a group.
\end{worked}

\begin{worked}{Question 3(a)}
If $[k]_{12}=[\ell]_{12}$, then $12\mid(k-\ell)$. Therefore $8\mid2(k-\ell)$, because $2(k-\ell)$ is a multiple of $24$. Hence $[2k]_8=[2\ell]_8$, proving well-definedness.

For any $[k]_{12},[\ell]_{12}$,
\[
\phi([k]_{12}+[\ell]_{12})=[2(k+\ell)]_8
=[2k]_8+[2\ell]_8
=\phi([k]_{12})+\phi([\ell]_{12}).
\]
Thus $\phi$ is a homomorphism. Its kernel consists of $[k]_{12}$ for which $[2k]_8=[0]_8$, meaning $8\mid2k$, or $4\mid k$. Hence
\[
\ker\phi=\{[0]_{12},[4]_{12},[8]_{12}\}.
\]
The outputs are exactly the even residue classes modulo $8$:
\[
\operatorname{im}\phi=\{[0]_8,[2]_8,[4]_8,[6]_8\}.
\]
\end{worked}

\begin{worked}{Question 3(b)}
In an additive group, $\angles{[d]}=\{[kd]:k\in\Z\}$. Since $\gcd(7,18)=1$, the least positive $k$ with $18\mid7k$ is $18$. Thus $[7]$ has order $18$ and generates all of $\Z/18\Z$:
\[
\angles{[7]}=\{[0],[1],\ldots,[17]\}.
\]
Since $\gcd(5,20)=5$, the least positive $k$ with $20\mid5k$ is $4$. The first four multiples are distinct, so
\[
\angles{[5]}=\{[0],[5],[10],[15]\}\subseteq\Z/20\Z.
\]
This uses the formula $\operatorname{ord}([d]_n)=n/\gcd(n,d)$.
\end{worked}

\begin{worked}{Question 4(a)}
The relation $yxyz^{-2}=e$ is equivalent to $yxy=z^2$, hence
\[
x=y^{-1}z^2y^{-1}.
\]
Thus $y,z$ generate $P$. Let $F(y,z)$ be the free group on two symbols $y,z$. The assignment sending its generators to the corresponding elements of $P$ extends to a surjective homomorphism $\alpha:F(y,z)\to P$.

In the other direction, assign the generators of $P$ the values
\[
y\mapsto y,\qquad z\mapsto z,\qquad x\mapsto y^{-1}z^2y^{-1}
\quad\text{in }F(y,z).
\]
The relator maps to
$y(y^{-1}z^2y^{-1})yz^{-2}=e$, so this assignment defines a homomorphism $\beta:P\to F(y,z)$. The composite $\beta\alpha$ fixes $y,z$, so it is the identity on $F(y,z)$. The composite $\alpha\beta$ fixes $y,z$ and also fixes $x$ by the displayed relation, so it is the identity on $P$. Therefore $P\cong F(y,z)$.
\end{worked}

\begin{worked}{Question 4(b)}
For $f,g\in G^X$, $f(x)g(x)\in G$ for every $x$, so the pointwise product is again a function $X\to G$ (closure). For $f,g,h\in G^X$ and arbitrary $x\in X$,
\[
((fg)h)(x)=(f(x)g(x))h(x)=f(x)(g(x)h(x))=(f(gh))(x),
\]
using associativity in $G$. Equality at every $x$ gives associativity of functions.

The identity is the constant function $E(x)=e_G$ because $(Ef)(x)=e_Gf(x)=f(x)=(fE)(x)$. For $f\in G^X$, define $f^{-1}(x)=f(x)^{-1}$. Then $(ff^{-1})(x)=e_G=(f^{-1}f)(x)$ for all $x$, so $f^{-1}$ is a two-sided inverse. Hence $G^X$ is a group.
\end{worked}

\begin{worked}{Question 5}
\begin{enumerate}[label=(\alph*)]
\item \textbf{True.} $x\in f^*(B\cup C)$ exactly when $f(x)\in B$ or $f(x)\in C$.
\item \textbf{False.} Let $X=\{1,2\}$, let $f$ be constant with value $a$, and take $A=\{1\}$, $D=\{2\}$. Then $f_*(A\cap D)=\varnothing$, while $f_*(A)\cap f_*(D)=\{a\}$.
\item \textbf{True.} If $y\notin f(X)$, its fibre $f^{-1}(\{y\})$ is empty; $\varnothing\subseteq A$ for every $A$.
\item \textbf{True.} $\gcd(3,16)=1$, so $[3]$ has order $16$ and generates all of $\Z/16\Z$.
\item \textbf{False.} The image of $[1]_7$ must have order dividing $7$. Since $\Z/15\Z$ has no nonzero element of order $7$, it must map to $[0]$, making the homomorphism trivial.
\item \textbf{True.} $g^2=e$ implies $g^{-1}=g$ for every $g$. Thus $ab=(ab)^{-1}=b^{-1}a^{-1}=ba$.
\end{enumerate}
\end{worked}

\end{document}
